\lstdefinestyle{Python}{
	language = {Python},
	style    = lgeneral,
}

% \subsubsection{包包包}

\begin{lstlisting}[style=Python]
# -*- coding: UTF-8 -*-
import sys, random, bisect                           # 二分查找
from collections import Counter, defaultdict, deque  # 常用集合类
from copy import deepcopy                            # 深拷贝
from functools import cmp_to_key, lru_cache, reduce  # 函数工具

# 堆操作
from heapq import merge, heapify, heappop, heappush, heappushpop, nlargest, nsmallest, heapreplace

# 迭代工具
from itertools import accumulate, combinations, permutations, count, product

# 运算符
from operator import add, iand, ior, itemgetter, mul, xor

# 字符串常量
from string import ascii_lowercase, ascii_uppercase, ascii_letters

# 类型提示
import typing
from typing import *

input = lambda: sys.stdin.readline().rstrip()

def solve() -> str:
    return "Hello World"

print('\n'.join(solve() for _ in range(int(input()))))
\end{lstlisting}

\newpage

% \subsubsection{递归深度}

使用 \texttt{sys} 里的函数来调整递归深度：

\begin{lstlisting}[style=Python]
sys.setrecursionlimit(int)   # 设置最大递归深度，这个是全局的
sys.getrecursionlimit()      # 当前允许的最大递归深度
\end{lstlisting}

% \subsubsection{排序}

使用 \texttt{cmp\_to\_key}，以从大到小排序为例：

\begin{lstlisting}[style=Python]
from functools import cmp_to_key
"""
-1: x 排在 y 前
0:  x 和 y 相同
1:  x 排在 y 后
"""
def cmp(x, y):
    if x == y:
        return 0
    return -1 if x > y else 1
a = list(map(int, input().split()))
a.sort(key=cmp_to_key(cmp))
b = sorted(a, key=cmp_to_key(cmp))
\end{lstlisting}

这个排序跟 C++ 不同，C++ 是：

\begin{itemize}
	\item \texttt{cmp(x, y) == 1}，则 $x$ 严格排在 $y$ 前
	\item \texttt{cmp(x, y) == 0}，则 $x$ 和 $y$ 相同或者严格排在 $y$ 后
\end{itemize}

此外，可以直接用 \texttt{a = sorted(map(int, input().split()))} 来直接得到\textbf{升序排序}后的 \texttt{list}

\newpage

% \subsubsection{日期与时间}

\begin{lstlisting}[style=Python]
from datetime import datetime, timedelta

d = datetime(1970, 1, 1, 1, 3, 5)
print(d.year, d.month, d.day)           # 1970 1 1
print(d.hour, d.minute, d.second)       # 1 3 5

# Thursday
print(d.weekday())                      # 3 in [0,6]
print(d.isoweekday())                   # 4 in [1,7]

year, week, weekday = d.isocalendar()   # 年份, 周数, 星期几
print(year, week, weekday)              # 1970 1 4

d = d.replace(year=2026, month=4)
print(d.strftime("%Y-%m-%d %H:%M:%S"))  # 2026-04-01 01:03:05
print(d.strftime("%A, %B %d, %Y"))      # Wednesday, April 01, 2026

dt = timedelta(days=10, hours=3)
d += dt
print(d)                                # 2026-04-11 04:03:05
print(dt.days, dt.total_seconds())      # 10 874800.0

try:
    d = datetime(2006, 2, 29)
    print(f"{d} is valid")
except ValueError:
    print(f"invalid date")
\end{lstlisting}
